\section{Related Work}
\label{sec:related_work}

\para{Truth and knowledge representations in \llms.}
A growing body of work investigates how \llms internally represent factual knowledge. \citet{marks2023geometry} discovered emergent linear structure in LLaMA representations of true/false statements, showing that truth value becomes linearly separable in middle layers and that this structure emerges with scale. \citet{burns2023discovering} proposed Contrast-Consistent Search (CCS) for unsupervised discovery of truth directions in activation space, finding that middle layers can outperform final layers for truth detection. \citet{azaria2023internal} demonstrated that classifiers trained on hidden states can detect when \llms generate falsehoods, while \citet{liu2023cognitive} examined the gap between internal truth representations and output probabilities. \citet{meng2022locating} used causal tracing to localize factual associations in specific MLP layers of GPT models. Our work extends this line of research by comparing truth representations with a distinct category---non-epistemic beliefs---that has not been probed before.

\para{Belief representations and theory of mind.}
\citet{zhu2024language} showed that \llms develop linearly decodable representations of belief states for both self and others, with oracle beliefs broadly encoded across layers 13--15 and protagonist beliefs concentrated in specific attention heads at layer 10 of Mistral-7B. \citet{bortoletto2024brittle} conducted the first systematic investigation of belief representations across 12 models (Pythia 70M--12B, LLaMA-2 7B--70B), finding that oracle beliefs peak in early layers while protagonist beliefs peak at intermediate layers and improve with fine-tuning. \citet{churina2025layer} tracked how belief representations shift under continual pretraining poisoning. Unlike these studies, which focus on belief \emph{states} (true belief vs.\ false belief in theory-of-mind scenarios), we probe belief \emph{types}---the distinction between factual knowledge claims and subjective opinions.

\para{Probing methodology.}
Linear probing was introduced by \citet{alain2016understanding} to study intermediate layer representations and has become a standard tool in interpretability research \citep{conneau2018probing,belinkov2022probing}. \citet{hewitt2019designing} proposed control tasks with random labels to measure selectivity---the gap between true-label and random-label probe accuracy---as a guard against probes that memorize rather than decode genuine structure. \citet{voita2020information} introduced information-theoretic probing with minimum description length for more faithful measurement. We adopt the mass-mean probing direction from \citet{marks2023geometry}, which is optimization-free and shown to be the most causally implicated probing direction, combined with the selectivity framework of \citet{hewitt2019designing}.

\para{Layer-wise processing in transformers.}
Research on transformer internal organization has revealed consistent patterns across architectures. \citet{tenney2019bert} showed that BERT represents classical NLP pipeline stages in sequence: POS tagging in lower layers, parsing in middle layers, and semantic roles in upper layers. \citet{jawahar2019does} confirmed that lower layers capture phrase-level syntax while upper layers capture semantic features. \citet{ju2024large} found that context knowledge is encoded increasingly in upper layers, with knowledge-related entity tokens encoding information earlier than other tokens. \citet{jin2024exploring} proposed the notion of ``concept depth,'' showing that more complex concepts are processed at deeper layers. Our findings complement this literature by characterizing the layer-wise profile of belief type and truth value encoding.

\para{Epistemic reasoning in \llms.}
\citet{suzgun2024belief} introduced the KaBLE benchmark with 13,000 questions testing \llms' ability to distinguish knowledge, belief, and fact. \citet{herrmann2024standards} established theoretical standards for what counts as genuine belief representation in \llms. \citet{vesga2025psychological} proposed a psychological framework distinguishing epistemic from non-epistemic beliefs, which directly motivates our probing categories. \citet{li2025representations} studied how \llms represent facts, fiction, and forecasts. Our work is the first to use this epistemic/non-epistemic distinction as the basis for a systematic probing study of internal representations.
